% CA23 — Full report (LaTeX)
\documentclass[11pt]{article}
\usepackage[utf8]{inputenc}
\usepackage{amsmath,amssymb,amsthm}
\usepackage{graphicx}
\usepackage{booktabs}
\usepackage{hyperref}
\usepackage{geometry}
\usepackage{caption}
\usepackage{subcaption}
\usepackage{algorithm}
\usepackage{algpseudocode}
\usepackage{enumitem}
\geometry{margin=1in}

\title{CA23: A Reproducible Scaffold for Policy-Gradient Experiments}
\author{Course Maintainers}
\date{\today}

\begin{document}

\maketitle

\begin{abstract}
This document provides a comprehensive, reproducible description of the CA23
scaffold: design goals, implementation details, algorithmic building blocks,
experimental protocols, and best practices for reproducible reinforcement
learning experiments. The report is intended as a short, self-contained piece
suitable for inclusion in a course submission or as the basis for a short
conference-style write-up. We include equations for the optimization
objectives, suggested hyperparameter sweeps, metrics and statistical tests,
and instructions for packaging and sharing results.
\end{abstract}

\section{Introduction}
Reinforcement learning (RL) experiments often require careful attention to
reproducibility: controlled randomness, configuration management, and clear
reporting of evaluation metrics. The CA23 scaffold provides a minimal,
well-documented toolkit for implementing policy-gradient style agents and
evaluating them on canonical control tasks. Our goals are (1) to make it easy
for students and researchers to implement and iterate on policy-gradient
algorithms, and (2) to ensure experiments are reproducible and reportable.

This report expands on the short `REPORT.md`: it formalizes the loss
functions, provides pseudocode, and specifies a complete experimental protocol
for a small set of baseline experiments (baseline, entropy sweep, network
capacity ablation, learning rate sensitivity).

\section{Related work}
Classic references include Sutton and Barto's textbook \cite{sutton_barto},
which derives policy-gradient estimators and advantage-based variance
reduction techniques. OpenAI Spinning Up provides practical implementation
advice and recipes for common algorithms (policy gradient, PPO) \cite{spinup}.
This scaffold intentionally targets the policy-gradient family and provides a
clear place to insert improvements such as GAE, baselines, or trust-region
updates.

\section{Methods}
\subsection{Policy and value parameterization}
We use simple multi-layer perceptrons for both policy and value networks. For a
categorical action space, the policy $\pi_\theta(a|s)$ is parameterized by a
network returning action logits; action probabilities are given by softmax.
The value network $V_\phi(s)$ returns a scalar estimate of expected return.

\subsection{Objectives}
Let $\tau = (s_0,a_0,r_0,\dots,s_T)$ be an episode. The return from time $t$ is
$$R_t = \sum_{k=t}^{T} \gamma^{k-t} r_k$$
The policy gradient objective can be written as
$$J(\theta) = \mathbb{E}_\tau\left[ \sum_{t=0}^{T} \log \pi_\theta(a_t|s_t) A_t \right],$$
where $A_t$ is an estimator of advantage (we use $A_t = R_t - V_\phi(s_t)$ as
a simple baseline). Minimizing the negative expected advantage gives the
loss used in code:
$$\mathcal{L}_\text{policy} = -\mathbb{E}_t[\log \pi_\theta(a_t|s_t) A_t].$$

We also include a value-function loss (MSE):
$$\mathcal{L}_\text{value} = \mathbb{E}_t\left[(V_\phi(s_t) - R_t)^2\right].$$
Entropy regularization encourages exploration by adding a term to encourage
high-entropy policies. Denoting the entropy of the policy at state $s$ by
$\mathcal{H}(\pi(\cdot|s))$, the combined loss minimized is
$$\mathcal{L} = \mathcal{L}_\text{policy} + c_v \mathcal{L}_\text{value} - \lambda \mathbb{E}_t[\mathcal{H}(\pi(\cdot|s_t))],$$
where $c_v$ is the value loss coefficient and $\lambda$ is the entropy
regularization coefficient (in our code `entropy_coef`).

\subsection{Numerical and implementation details}
- Policy action selection uses PyTorch's `Categorical` distribution to obtain
  actions and log-probabilities. For deterministic evaluation we pick the
  argmax of the probabilities.
- Careful dtype handling is used in `src/losses.py` so tests remain stable with
  mixed dtypes.
- To support multiple gym versions, we handle both the older `(obs, reward,
  done, info)` and newer `(obs, reward, terminated, truncated, info)` step
  signatures.

\section{Implementation files and responsibilities}
Table~\ref{tab:filemap} lists the principal files and their responsibilities.

\begin{table}[h]
\centering
\begin{tabular}{@{}ll@{}}\toprule
File & Purpose \\\midrule
`src/config.py` & `ExperimentConfig` dataclass: YAML load/save, validation \\
`src/model.py` & `PolicyNetwork`, `ValueNetwork` (MLP-based) \\
`src/losses.py` & `policy_gradient_loss`, `value_loss`, `entropy_loss` (compat wrapper) \\
`src/data.py` & `collect_episodes` (lazy gym import), `discounts` \\
`src/utils.py` & `set_seed`, `get_device`, `save_figure`, `ensure_dir` \\
`scripts/train.py` & Pedagogical training loop and CSV logging \\
`tests/` & Unit tests and import-safe smoke checks \\
\bottomrule
\end{tabular}
\caption{File map and responsibilities for the CA23 scaffold.}
\label{tab:filemap}
\end{table}

\section{Experimental design}
We provide a canonical experimental protocol intended to be reproducible and
interpretable. Experiments recommended for a short course project include:

\begin{enumerate}[itemsep=2pt]
\item Baseline: run the default config on \texttt{CartPole-v1} for 5 seeds.
\item Entropy sweep: \(\lambda \in \{0.0, 0.01, 0.1\}\) across 5 seeds.
\item Network capacity ablation: hidden sizes in \{(32,32), (64,64), (128,128)\}.
\item LR sensitivity: learning rates \{1e-4, 1e-3, 3e-3\}.
\end{enumerate}

For each experiment report the per-episode return, policy and value losses,
entropy, and the wall-clock time per episode. Use a consistent stopping
criterion or fixed number of episodes across runs for fair comparison.

\subsection{Hyperparameters}
Table~\ref{tab:hparams} gives suggested default hyperparameters used in the
scaffold and in the baseline experiments.

\begin{table}[h]
\centering
\begin{tabular}{@{}ll@{}}\toprule
Hyperparameter & Value \\\midrule
env_name & CartPole-v1 \\
seed & 0 (sweep seeds 0–4) \\
learning_rate & 1e-3 \\
gamma & 0.99 \\
hidden_sizes & (64,64) \\
batch_size & 32 \\
max_episodes & 200 \\
entropy_coef & 0.0 (sweep as above) \\
\bottomrule
\end{tabular}
\caption{Suggested default hyperparameters.}
\label{tab:hparams}
\end{table}

\section{Evaluation metrics and statistical tests}
Primary metric: mean episode return computed per episode and smoothed with a
moving average (window 50). For comparing conditions, compute the mean and
standard deviation of final returns across seeds and report confidence
intervals. For pairwise comparisons use a two-sided t-test and report p-value
and Cohen's d for effect size.

\section{Example results (placeholder)}
To include real results, run the sweep above and place aggregated CSVs under
`runs/`. For this template we include placeholders shown in
Figure~\ref{fig:placeholder} and Table~\ref{tab:results}.

\begin{figure}[h]
\centering
\fbox{\parbox[c][0.25\textheight][c]{0.6\textwidth}{\centering Placeholder figure: run experiments to generate plots}}
\caption{Placeholder plot location — replace with learning curve generated from `runs/`.}
\label{fig:placeholder}
\end{figure}

\begin{table}[h]
\centering
\begin{tabular}{@{}lcc@{}}\toprule
Experiment & Mean return & Std \\\midrule
Baseline & -- & -- \\
Entropy 0.01 & -- & -- \\
Entropy 0.1 & -- & -- \\
\bottomrule
\end{tabular}
\caption{Placeholder results table (fill in after running experiments).}
\label{tab:results}
\end{table}

\section{Reproducibility checklist}
Follow this checklist to ensure reproducible runs and publishable artifacts:
\begin{itemize}
\item Save `config.yaml` for each run under `runs/<id>/config.yaml` using
  `ExperimentConfig.to_yaml`.
\item Record git commit hash (e.g., `git rev-parse HEAD`) and include it in the
  run metadata.
\item Store `pip freeze` in the run folder or an environment YAML.
\item Record seeds and, when using GPUs, device identifiers.
\item Save per-episode logs and model checkpoints.
\item Report aggregated results with mean ± std across seeds and include the
  raw data (CSV) in a `results/` folder for easy re-analysis.
\end{itemize}

\section{Appendix A: Algorithm pseudocode}
\begin{algorithm}[h]
\caption{Vanilla Policy Gradient (episodic)}
\begin{algorithmic}[1]
\State Initialize policy $\theta$ and value $\phi$.
\For{episode = 1 to N}
  \State Collect trajectory $\tau$ by rolling out $\pi_\theta$.
  \State Compute returns $R_t$ for $t$ in trajectory.
  \State Compute advantages $A_t = R_t - V_\phi(s_t)$.
  \State $\theta \leftarrow \theta - \alpha \nabla_\theta [ -\mathbb{E}_t[\log \pi_\theta(a_t|s_t) A_t] - \lambda \mathbb{E}_t[\mathcal{H}(\pi(\cdot|s_t)) ] ]$
  \State $\phi \leftarrow \phi - \beta \nabla_\phi \mathbb{E}_t[(V_\phi(s_t) - R_t)^2]$
\EndFor
\end{algorithmic}
\end{algorithm}

\section{Appendix B: How to run}
A minimal sequence of commands to run the baseline experiment:
\begin{verbatim}
python -m venv .venv && source .venv/bin/activate
python -m pip install -r requirements.txt
python scripts/train.py --config configs/debug.yaml
\end{verbatim}

To run a sweep, create per-run YAML configs in `runs/<exp>/<seed>/config.yaml`
and launch `scripts/train.py` with those configs (see the notebook for an
example shell loop).

\section{Appendix C: API reference (short)}
\begin{itemize}
\item `src.config.ExperimentConfig` — dataclass with fields shown in Table~\ref{tab:hparams}. Use `from_yaml` and `to_yaml` for load/save.
\item `src.model.PolicyNetwork` — constructor `(obs_dim, action_dim, hidden_sizes=(64,64))`. Use `.get_action(obs_tensor)` to sample actions.
\item `src.losses.policy_gradient_loss` — signature `(log_probs, advantages)`.
\item `src.data.collect_episodes` — `(env_name, policy, num_episodes=1, max_steps=1000)`; returns list of episodes (each a list of `Transition`).
\end{itemize}

\section{Conclusion}
We present CA23 as a minimal but practical scaffold for policy-gradient
experiments with an emphasis on reproducibility and reporting. The included
code, tests, and documentation provide a clear starting point for course
projects and small-scale research experiments. We encourage users to extend the
scaffold with more advanced algorithms and improved logging/checkpointing for
larger studies.

\section*{Acknowledgements}
Thanks to the course staff and maintainers for continuous improvement of the
scaffold.

\begin{thebibliography}{9}
\bibitem{sutton_barto}
Richard S. Sutton and Andrew G. Barto. Reinforcement Learning: An Introduction. 2nd edition, MIT Press, 2018.

\bibitem{spinup}
OpenAI. Spinning Up in Deep RL. 2018. \url{https://spinningup.openai.com}
\end{thebibliography}

\end{document}
